
\begin{table}[t!]
\centering
\caption{Baseline benchmark under original and selected augmented feature sets. Entries report the mean test set Threat Score across five zone-stratified split seeds.}
\label{tab:baseline_original_augmented_features}
\small
\setlength{\tabcolsep}{2pt}
\begin{tabularx}{\linewidth}{@{}Xcccc@{}}
\toprule
\textbf{Method} & \textbf{50/25/25} & \textbf{60/20/20} & \textbf{70/15/15} & \textbf{80/10/10} \\
\midrule
USGS~\cite{staley2016} 
& $0.388 \pm 0.031$ 
& $0.371 \pm 0.030$ 
& $0.376 \pm 0.030$ 
& $0.402 \pm 0.044$ \\
\midrule
LR~\cite{nikolopoulos2018pfdf_rf, kern2017ml_pfdf} 
& $0.449 \pm 0.022$ 
& $0.409 \pm 0.052$ 
& $0.441 \pm 0.035$ 
& $0.484 \pm 0.035$ \\
CART~\cite{kern2017ml_pfdf} 
& $0.523 \pm 0.032$ 
& $0.509 \pm 0.043$ 
& $0.556 \pm 0.078$ 
& $0.529 \pm 0.042$ \\
SVM~\cite{kern2017ml_pfdf} 
& $0.504 \pm 0.046$ 
& $0.464 \pm 0.034$ 
& $0.524 \pm 0.058$ 
& $0.552 \pm 0.066$ \\
RF~\cite{nikolopoulos2018pfdf_rf} 
& \underline{$0.585 \pm 0.027$} 
& $\bm{0.617 \pm 0.037}$ 
& $\bm{0.658 \pm 0.046}$ 
& $\bm{0.674 \pm 0.041}$ \\
HGB 
& $\bm{0.604 \pm 0.038}$ 
& \underline{$0.611 \pm 0.032$} 
& \underline{$0.640 \pm 0.040$} 
& \underline{$0.648 \pm 0.052$} \\
NN~\cite{kern2017ml_pfdf} 
& $0.540 \pm 0.030$ 
& $0.558 \pm 0.046$ 
& $0.583 \pm 0.032$ 
& $0.613 \pm 0.012$ \\
\midrule
LR$^{*}$ 
& $0.457 \pm 0.025$ 
& $0.483 \pm 0.022$ 
& $0.481 \pm 0.021$ 
& $0.483 \pm 0.033$ \\
CART$^{*}$ 
& $0.504 \pm 0.035$ 
& $0.542 \pm 0.044$ 
& $0.542 \pm 0.050$ 
& $0.549 \pm 0.030$ \\
SVM$^{*}$ 
& $0.547 \pm 0.020$ 
& \underline{$0.597 \pm 0.029$} 
& \underline{$0.648 \pm 0.025$} 
& \underline{$0.670 \pm 0.033$} \\
RF$^{*}$ 
& \underline{$0.600 \pm 0.037$} 
& $\bm{0.634 \pm 0.046}$ 
& $\bm{0.685 \pm 0.048}$ 
& $\bm{0.687 \pm 0.065}$ \\
HGB$^{*}$ 
& $0.595 \pm 0.062$ 
& $0.605 \pm 0.031$ 
& $0.634 \pm 0.031$ 
& $0.652 \pm 0.028$ \\
NN$^{*}$ 
& $\bm{0.604 \pm 0.021}$ 
& \underline{$0.611 \pm 0.041$} 
& $0.632 \pm 0.018$ 
& $0.663 \pm 0.043$ \\
\bottomrule
\end{tabularx}
\end{table}

\section{Experimental Results}

\subsection{Benchmarking Performance}

Table~\ref{tab:baseline_original_augmented_features} compares classical and neural baselines under both the original and the selected augmented feature table. The fixed USGS logistic-regression (USGS) model is included as an operational reference and is not retrained. All other methods follow the same repeated zone-stratified evaluation protocol, with entries reporting the mean test Threat Score and standard deviation across five split seeds. While the original feature table already provides strong performance for tree-based ensembles, the selected augmented features further improve most model families, particularly SVM, random forest, and neural network models. The widely deployed USGS logistic-regression baseline performs substantially worse than modern learning methods. This is primarily because its fixed coefficients cannot adaptively capture the diverse regional mechanisms underlying PFDF formation, an issue we analyze in Section~\ref{sec-analysis}.

\begin{figure}[t!]
     \centering
     \includegraphics[width=1\linewidth]{Figure/top10.pdf}
     \caption{Top-10 Important Features for PFDF prediction across regions with color denoting feature category.}
\Description{Top-10 Important Features for PFDF prediction across regions with color denoting feature category.}
     \label{fig:feature_importance}
\end{figure}

\subsection{Global Feature Importance}
After applying the RL feature selection in Section~\ref{sec-RL}, we visualize the top-10 important features across different zones

From the global feature importance analysis, storm-related variables dominate the most influential features, aligning with prior literature identifying storm conditions as the primary triggering factor for PFDF events~\cite{staley2016}. Meanwhile, soil erodibility, burn severity, and terrain slope also consistently emerge as critical contributing factors. \textit{Notably, the geographical zone itself appears as a highly important feature, highlighting the strong region-specific characteristics underlying PFDF formation.} This observation suggests that different regions exhibit distinct interactions among meteorological, wildfire, soil, and terrain conditions, further motivating the need for adaptive and region-aware prediction models. This finding also resonates with substantially lower performance of the rigid fixed-coefficient USGS method under the zone-stratified evaluation setting in Table~\ref{tab:baseline_original_augmented_features}, as a single globally calibrated logistic formulation may struggle to capture region-dependent PFDF mechanisms.

Comparing different zones, we can clearly observe substantial variation in feature importance across regions. OSD and SGSBSJ, located in California, are primarily dominated by storm/rainfall-related factors, indicating the strong influence of meteorological triggering conditions in these regions. In contrast, for SWCO and WMT, terrain-related and geographical features emerge as the most influential factors, suggesting that topographical and regional characteristics play a more dominant role in determining PFDF occurrence in these areas.

\begin{table}[t!]
\centering
\caption{Region-wise absolute Threat Score improvement of the original-feature neural network over the fixed USGS logistic-regression reference. Entries report $\Delta\mathrm{TS}=\mathrm{TS}_{\mathrm{NN}}-\mathrm{TS}_{\mathrm{USGS}}$ averaged across split seeds. Region counts are reported as total records (positive/negative).}
\label{tab:region_nn_usgs_improvement}
\small
\setlength{\tabcolsep}{1.2pt}
\begin{tabular}{l|llcccc}
\toprule
\textbf{Area} & \textbf{Zone} & \textbf{Records} & \textbf{50/25/25} & \textbf{60/20/20} & \textbf{70/15/15} & \textbf{80/10/10} \\
\midrule
\multirow{3}{*}{\makecell[l]{\textbf{Southern} \\\textbf{California}}}
& OSD    & 103 (25/78)   & \cellcolor{ImproveOrange}$+0.19$ & \cellcolor{ImproveOrange}$+0.23$ & \cellcolor{ImproveOrange}$+0.22$ & \cellcolor{ImproveOrange}$+0.19$ \\
& SGSBSJ & 646 (176/470) & \cellcolor{ImproveOrange}$+0.16$ & \cellcolor{ImproveOrange}$+0.19$ & \cellcolor{ImproveOrange}$+0.19$ & \cellcolor{ImproveOrange}$+0.17$ \\
& VEN    & 24 (0/24)     & \cellcolor{NoChangeGray}$+0.00$ & \cellcolor{NoChangeGray}$+0.00$ & \cellcolor{NoChangeGray}N/A      & \cellcolor{NoChangeGray}$+0.00$ \\
\midrule
\multirow{2}{*}{\textbf{Arizona}}
& NAZ    & 71 (5/66)     & \cellcolor{NoChangeGray}$+0.05$ & \cellcolor{NoChangeGray}$+0.12$ & \cellcolor{NoChangeGray}$+0.00$ & \cellcolor{NoChangeGray}$+0.00$ \\
& SAZ    & 13 (2/11)     & \cellcolor{NoChangeGray}$+0.00$ & \cellcolor{NoChangeGray}$+0.00$ & \cellcolor{NoChangeGray}$+0.00$ & \cellcolor{NoChangeGray}$+0.00$ \\
\midrule
\textbf{N-Mexico}
& CNM    & 52 (35/17)    & \cellcolor{NoChangeGray}$-0.01$ & \cellcolor{NoChangeGray}$+0.00$ & \cellcolor{NoChangeGray}$+0.05$ & \cellcolor{NoChangeGray}$+0.06$ \\
\midrule
\multirow{3}{*}{\makecell[l]{\textbf{Colorado}\\\textbf{Western}\\\textbf{Plateau}}}
& FRCO   & 22 (4/18)     & \cellcolor{NoChangeGray}$-0.07$ & \cellcolor{NoChangeGray}$-0.20$ & \cellcolor{NoChangeGray}$+0.00$ & \cellcolor{NoChangeGray}$+0.00$ \\
& SWCO   & 15 (10/5)     & \cellcolor{ImproveOrange}$+0.30$ & \cellcolor{ImproveOrange}$+0.33$ & \cellcolor{ImproveOrange}$+0.50$ & \cellcolor{ImproveOrange}$+0.52$ \\
& WCO    & 145 (17/128)  & \cellcolor{ImproveOrange}$+0.16$ & \cellcolor{ImproveOrange}$+0.25$ & \cellcolor{ImproveOrange}$+0.15$ & \cellcolor{ImproveOrange}$+0.24$ \\
\midrule
\makecell[l]{\textbf{Northern}\\\textbf{Rockies}}
& MT     & 40 (18/22)    & \cellcolor{DecreaseBlue}$-0.19$ & \cellcolor{DecreaseBlue}$-0.17$ & \cellcolor{DecreaseBlue}$-0.14$ & \cellcolor{DecreaseBlue}$-0.18$ \\
\bottomrule
\end{tabular}
\end{table}
\subsection{Further Analysis}\label{sec-analysis}
We further diagnose why the neural model improves over the fixed USGS logistic-regression reference in some regions but not others. Figure~\ref{fig:regional_usgs_diagnostic} compares three representative zones. In each row, the left panel reports the ratio between the PFDF-event mean and the non-PFDF-event mean for the USGS inputs ($R$, $X_1$, $X_2$, $X_3$) and high-ranked learned features from Figure~\ref{fig:feature_importance}; values above one indicate larger feature values in observed PFDF events. The right panel shows the fixed USGS probability distribution by observed class, together with the 0.5 decision threshold and the missed PFDF count. 
% This setting separates two questions: whether the fixed USGS variables already align with local PFDF occurrence, and whether additional learned features expose stronger regional separation.


\subsubsection{OSD/SGSBSJ} both belong to the broader Southern California and show similar improvement patterns. Figure~\ref{fig:regional_usgs_diagnostic} reports \texttt{OSD} as the representative case. In \texttt{OSD}, the neural model improves substantially over the fixed-coefficient USGS baseline because several non-USGS or expanded storm descriptors separate PFDF from non-PFDF events more strongly than the original formula terms. For example, storm accumulation, 30/60-minute rain, 15-minute intensity, and average storm intensity all show large PFDF/non-PFDF ratios, exceeding several of the original USGS inputs. The fixed formula, therefore, misses 9 of 25 PFDF events, leaving room for the neural model to learn additional storm-response structure.
% beyond the four hand-designed USGS terms.

\subsubsection{SWCO} is the clearest failure mode for the fixed USGS rule. The USGS equation assumes that larger rainfall forcing should monotonically increase PFDF probability~\cite{staley2016, staley2016updated, staley2017prediction}, but in this zone, the observed PFDF events have a smaller mean $R$ than non-PFDF events, violating this monotonic trend. This reversed local relationship causes the fixed model to rank many PFDF events too low, missing 8 of 10 positives. The neural model can partially correct this by using the zone representation and other covariates to learn a region-specific relationship, including cases where rainfall-related effects are weaker or negative relative to the global rule. This interpretation is consistent with Figure~\ref{fig:feature_importance}, where geographic zone is highly ranked for \texttt{SWCO}, suggesting that location helps the model adapt the mechanism rather than applying one universal coefficient pattern. This is consistent with distinctive PFDF triggering conditions deduced from physics-based models. 

\subsubsection{MT} illustrates the opposite regime. Here, the fixed USGS variables already align well with the observed outcome: $R$, $X_1$, and $X_2$ are larger for PFDF events than for non-PFDF events, while $X_3$ is close to parity. As a result, the fixed model misses only 1 of 18 PFDF events. Because the baseline already captures most positives in this region, there is less capacity for a flexible neural model to improve, and additional learned interactions can slightly reduce performance under the repeated split setting.


\begin{figure}[t!]
\centering
\includegraphics[width=\linewidth]{Figure/USGS_analysis.pdf}
\caption{Regional diagnostic for OSD, SWCO, and MT. Left: PFDF/non-PFDF mean ratios for USGS inputs and high-ranked learned features. Right: fixed USGS predicted probability by observed class, with the decision threshold and missed PFDF events.}
\Description{Regional diagnostic for OSD, SWCO, and MT. Left: PFDF/non-PFDF mean ratios for USGS inputs and high-ranked learned features. Right: fixed USGS predicted probability by observed class, with the decision threshold and missed PFDF events.}
\label{fig:regional_usgs_diagnostic}
\end{figure}

Overall, Figure~\ref{fig:regional_usgs_diagnostic} shows that regional PFDF mechanisms are not uniform. Some regions benefit from additional storm and susceptibility features, some require an adaptive correction to the sign or strength of the rainfall relationship, and some are already well described by the fixed USGS structure. This supports the need for models that empirically learn region-dependent mechanisms.
% In this section, we present a physics-informed neural framework for PFDF prediction with two complementary components. First, we introduce a physics-guided knowledge augmentation module that integrates additional heterogeneous information, including meteorological and geomorphological factors. Second, we develop a physics-guided architecture that learns a structured dependency graph over feature variables, enabling physically consistent interactions while improving interpretability and predictive reliability.

% \subsection{Physics-guided Knowledge Augmentation}
% In addition to the base structured features, we use physics-guided knowledge augmentation to incorporate two complementary external information sources: weather context and soil--terrain susceptibility. The weather features describe the pre-event meteorological setting that controls rainfall triggering, while the soil and terrain features describe the watershed conditions that control runoff generation, sediment mobility, and hillslope response. We next describe how these additional features are curated and encoded.
% \subsubsection{Weather}
% Weather augmentation is motivated by the fact that PFDF occurrence depends not only on observed rain-gauge summaries, but also on the larger storm environment. Forecast discussions may indicate whether rainfall is expected to be convective, short-duration, spatially localized, persistent over several hours, terrain-enhanced, or explicitly hazardous for burn scars. These factors are physically relevant because PFDF initiation is sensitive to rainfall intensity, storm duration, antecedent wetness, orographic enhancement, and flash-flood or debris-flow concern.

% For each storm--watershed event, we query National Weather Service Area Forecast Discussions issued during the 48-hour window before storm onset and attach the most recent discussion issued at or before the storm start as pre-response forecast context. We consider two weather representations. The first, \textbf{NN+W}, appends a dense embedding of the full forecast narrative to the base structured features. The second, \textbf{NN+SW}, converts the same forecast text into 18 numeric structured weather indicators, as illustrated in Table~\ref{tab:prompt_weather_example}, including precipitation strength, short-duration intensity, storm duration, localized extreme rainfall, repeated or training cells, antecedent wetness, convection, orographic enhancement, flash-flood language, burn-scar sensitivity, monsoon or tropical moisture, slow storm motion, and an overall PFDF weather-risk score. The dense representation preserves broad language semantics, while the structured representation provides compact and physically interpretable forecast variables.


% \begin{table}[t]
% \centering
% \caption{Main performance benchmark under repeated zone-stratified splits. Entries are mean test Threat Score $\pm$ standard deviation.}
% \label{tab:main_benchmark}
% \small
% \setlength{\tabcolsep}{4pt}
% \resizebox{\linewidth}{!}{%
% \begin{tabular}{lcccc}
% \toprule
% \textbf{Model} & \textbf{50/25/25} & \textbf{60/20/20} & \textbf{70/15/15} & \textbf{80/10/10} \\
% \midrule
% P-LR & $0.388 \pm 0.031$ & $0.371 \pm 0.030$ & $0.376 \pm 0.030$ & $0.402 \pm 0.044$ \\
% NN & $0.540 \pm 0.030$ & $0.558 \pm 0.046$ & $0.583 \pm 0.032$ & $0.613 \pm 0.012$ \\
% NN+W & $0.600 \pm 0.030$ & $0.604 \pm 0.045$ & $0.622 \pm 0.025$ & $0.615 \pm 0.026$ \\
% NN+SW & $0.584 \pm 0.021$ & $0.604 \pm 0.041$ & $0.623 \pm 0.029$ & $0.639 \pm 0.028$ \\
% NN+AUG & $0.562 \pm 0.040$ & $0.581 \pm 0.046$ & $0.592 \pm 0.025$ & $0.636 \pm 0.058$ \\
% NN+selected SW+AUG & $0.604 \pm 0.021$ & $0.611 \pm 0.041$ & $0.632 \pm 0.018$ & $0.662 \pm 0.043$ \\
% % NN+G & $0.590 \pm 0.029$ & $0.591 \pm 0.026$ & $0.617 \pm 0.025$ & $0.648 \pm 0.045$ \\
% \bottomrule
% \end{tabular}
% }
% \end{table}

% \begin{table}[t]
% \centering
% \caption{Paired statistical tests over the 20 matched split-seed settings. Positive $\Delta$TS indicates improvement of the first model over the second.}
% \label{tab:stat_tests}
% \small
% \setlength{\tabcolsep}{4pt}
% \resizebox{\linewidth}{!}{%
% \begin{tabular}{lcccc}
% \toprule
% \textbf{Comparison} & \textbf{$n$} & \textbf{Mean $\Delta$TS} & \textbf{95\% bootstrap CI} & \textbf{Paired $t$ p} \\
% \midrule
% NN $-$ P-LR & 20 & $+0.1892$ & $[+0.1663,+0.2129]$ & $1.66{\times}10^{-12}$ \\
% NN+W $-$ NN & 20 & $+0.0365$ & $[+0.0201,+0.0514]$ & $1.32{\times}10^{-4}$ \\
% NN+AUG $-$ NN & 20 & $+0.0192$ & $[-0.0004,+0.0368]$ & $3.27{\times}10^{-2}$ \\
% \bottomrule
% \end{tabular}
% }
% \end{table}

% \subsubsection{Soil and Terrain}
% Soil and terrain augmentation is motivated by the distinction between rainfall triggering and watershed susceptibility. Antecedent soil moisture affects infiltration capacity and runoff production before a storm begins. Wetter pre-storm conditions can lower the additional rainfall needed for rapid runoff, while dry conditions may change the balance between initial abstraction, post-fire water repellency, and sediment entrainment. We retrieve antecedent soil moisture from NLDAS-2 Noah products using the depth-integrated \texttt{SoilM\_0\_100cm} variable in kg m$^{-2}$~\cite{xia2012nldas}. In the model tables, \texttt{soilm\_7d} denotes the calendar month immediately before storm start and is therefore a short antecedent proxy, while \texttt{soilm\_30d} denotes the mean of the prior three calendar months and represents longer-term wetness.

% Terrain variables describe the static geomorphic setting that shapes how rainfall is converted into runoff and sediment transport. We summarize each watershed segment using mean slope and mean aspect. Slope captures gravitational driving stress, flow concentration potential, and the likelihood that runoff can mobilize sediment on steep burned hillslopes. Aspect provides a compact proxy for radiation exposure, drying, snow persistence, vegetation recovery, and moisture-retention differences. Together, antecedent soil moisture, slope, and aspect extend the original rainfall and burn-severity table toward a more complete storm--landscape state: rainfall describes the trigger, while soil and terrain describe the watershed's readiness to respond.


% \newpage
% \subsection{Physics-guided Neural Architecture}

% Feature augmentation adds information to the input vector, but it does not control how the model combines variables. We therefore also study architecture augmentation through a graph-constrained neural model, denoted \textbf{NN+G}. The goal is to impose a screened physical dependency structure on neural prediction so that interactions among storm, landscape, soil-moisture, terrain, season, zone, and PFDF nodes are disciplined by a graph rather than learned freely by a fully connected MLP.

% \subsubsection{Graph Modeling: Nodes and Edges}

% Each scalar predictor is treated as a graph node. The zone variable is represented by a learned embedding node rather than as a continuous scalar. Edges encode candidate directional dependencies among predictors and the PFDF response. For example, burn and terrain variables can affect runoff sensitivity, rainfall variables can affect triggering intensity, antecedent soil moisture can influence runoff production, and season or zone can influence storm and landscape context. The graph is used as a dependency prior for prediction, not as a claim that every retained edge is definitively causal.

% \subsubsection{Edge Construction and Modification}

% The dependency graph combines domain priors with data-driven screening. Candidate edges are informed by PFDF process knowledge, then refined using PC-style conditional-independence discovery~\cite{spirtes2000}, DirectLiNGAM~\cite{shimizu2006}, and downstream edge validation. We use this hybrid construction because purely learned graphs can be unstable on small environmental datasets, while purely hand-specified graphs may miss empirical dependencies. The final graph is therefore best interpreted as a physically screened dependency graph.

% \subsubsection{Graph-Based MLP Prediction}

% NN+G embeds each node into a hidden representation, propagates messages only along allowed graph edges, and maps the resulting node states to a debris-flow logit through a readout MLP. If $A$ denotes the graph adjacency matrix and $\mathbf{h}^{(0)}$ denotes the initial node embeddings, a graph layer updates each node using only its graph parents:
% \[
% \mathbf{h}^{(1)}_j = \phi(\mathbf{W}_s \mathbf{h}^{(0)}_j +
% \sum_{i:A_{ij}=1} \mathbf{W}_m \mathbf{h}^{(0)}_i),
% \]
% where $\phi$ is a nonlinear activation. This architecture uses the same physical input information as NN+AUG, allowing the experiment to isolate whether graph-constrained information flow improves prediction beyond ordinary feature augmentation.
